% Part: incompleteness
% Chapter: incompleteness-provability
% Section: second-incompleteness-thm

\documentclass[../../../include/open-logic-section]{subfiles}

\begin{document}

\olfileid{inc}{inp}{2in}

\olsection{قضیهٔ دوم ناتمامیت}

چگونه می‌توانیم این ادعا را بیان کنیم که~$\Th{PA}$ سازگاری خودش را
اثبات نمی‌کند؟ ناسازگاری~$\Th{PA}$ همان است که
$\Th{PA} \Proves \eq[0][1]$. پس می‌توانیم گزارهٔ سازگاری
$\OCon[\Th{PA}]$ را !!{sentence}~$\lnot
\OProv[\Th{PA}](\gn{\eq[0][1]})$ بگیریم؛ آنگاه قضیهٔ زیر مقصود را
برآورده می‌کند:

\begin{thm}
\ollabel{thm:second-incompleteness}
با فرض سازگاری~$\Th{PA}$، نظریهٔ~$\Th{PA}$، $\OCon[\Th{PA}]$ را
!!{derive} نمی‌کند.
\end{thm}

مهم است توجه کنیم که قضیه به بازنمایی خاص~$\OCon[\Th{PA}]$ (یعنی
بازنمایی خاص~$\OProv[\Th{PA}](y)$) وابسته است. تنها چیزی که به کار
خواهیم برد این است که بازنمایی~$\OProv[\Th{PA}](y)$ سه شرط
!!{derivability} را برآورده می‌کند؛ پس قضیه به هر نظریه‌ای تعمیم می‌یابد
که !!a{derivability} محمول دارای این ویژگی‌ها داشته باشد.

خواندن طرح استدلال گودل آموزنده است، زیرا قضیه همچون نتیجه‌گیری‌ای
درخشان از آن حاصل می‌شود. استدلال چنین است. بگذارید~$!G_\Th{PA}$ جملهٔ
گودلی باشد که در اثبات~\olref[1in]{thm:first-incompleteness} ساختیم.
نشان داده‌ایم: «اگر~$\Th{PA}$ سازگار باشد، آنگاه~$\Th{PA}$،
$!G_\Th{PA}$ را !!{derive} نمی‌کند.» اگر این را \emph{درون}
$\Th{PA}$ صوری کنیم، اثباتی از عبارت زیر داریم:
\[
\OCon[\Th{PA}] \lif \lnot \OProv[\Th{PA}](\gn{!G_\Th{PA}}).
\]
اکنون فرض کنید~$\Th{PA}$، $\OCon[\Th{PA}]$ را !!{derive}s می‌کند.
آنگاه~$\lnot \OProv[\Th{PA}](\gn{!G_\Th{PA}})$ را !!{derive}s می‌کند.
اما چون~$!G_\Th{PA}$ یک جملهٔ گودلی است، این عبارت با~$!G_\Th{PA}$
هم‌ارز است. پس~$\Th{PA}$، $!G_\Th{PA}$ را !!{derive}s می‌کند.

اما می‌دانیم اگر~$\Th{PA}$ سازگار باشد، $!G_\Th{PA}$ را !!{derive}
نمی‌کند!{} پس اگر~$\Th{PA}$ سازگار باشد، نمی‌تواند
$\OCon[\Th{PA}]$ را !!{derive} کند.

برای دقیق‌ترکردن استدلال، $!G_\Th{PA}$ را جملهٔ گودلیِ~$\Th{PA}$
می‌گیریم و با استفاده از شرایط !!{derivability} یعنی (P1)--(P3) نشان
می‌دهیم~$\Th{PA}$، $\OCon[\Th{PA}] \lif !G_\Th{PA}$ را !!{derive}s
می‌کند. این نشان خواهد داد که~$\Th{PA}$، $\OCon[\Th{PA}]$ را
!!{derive} نمی‌کند. در ادامه طرح اثبات در~$\Th{PA}$ آمده است. (برای
سادگی، زیرنویس‌های~$\Th{PA}$ را حذف می‌کنیم.)
\begin{align}
& !G \liff \lnot \OProv(\gn{!G}) \ollabel{G2-1}\\
& \qquad\text{$!G$ یک جملهٔ گودلی است}\notag \\
& !G \lif \lnot \OProv(\gn{!G}) \ollabel{G2-2}\\
  & \qquad\text{از \olref{G2-1}} \notag\\
& !G \lif
  (\OProv(\gn{!G}) \lif \lfalse) \ollabel{G2-3}\\
  & \qquad\text{از \olref{G2-2} به‌کمک منطق}\notag\\
& \OProv(\gn{
    !G \lif
    (\OProv(\gn{!G}) \lif \lfalse)
  }) \ollabel{G2-4}\\
  & \qquad\text{از \olref{G2-3} بنا بر شرط P1} \notag\\
& \OProv(\gn{!G}) \lif
  \OProv(\gn{
    (\OProv(\gn{!G}) \lif \lfalse)
    }) \ollabel{G2-5}\\
  & \qquad\text{از \olref{G2-4} بنا بر شرط P2} \notag\\
& \OProv(\gn{!G}) \lif (\OProv(\gn{\OProv(\gn{!G})}) \lif \OProv(\gn{\lfalse})) \ollabel{G2-6}\\
  & \qquad\text{از \olref{G2-5} بنا بر شرط P2 و منطق} \notag\\
& \OProv(\gn{!G}) \lif
  \OProv(\gn{\OProv(\gn{!G})}) \ollabel{G2-7}\\
   & \qquad\text{بنا بر P3} \notag\\
& \OProv(\gn{!G}) \lif \OProv(\gn{\lfalse}) \ollabel{G2-8}\\
  & \qquad \text{از \olref{G2-6} و \olref{G2-7} به‌کمک منطق}\notag\\
& \OCon \lif \lnot \OProv(\gn{!G}) \ollabel{G2-9}\\
  & \qquad\text{عکس نقیض \olref{G2-8} و $\OCon \ident \lnot \OProv(\gn{\lfalse})$}\notag \\
& \OCon \lif !G \notag\\
  & \qquad\text{از \olref{G2-1} و \olref{G2-9} به‌کمک منطق}\notag
\end{align}
کاربرد منطق در بالا فقط از واقعیت‌های مقدماتی منطق گزاره‌ای بهره
می‌گیرد؛ برای نمونه، \olref{G2-3} از
$\Proves \lnot!A \liff (!A\lif \lfalse)$ و~\olref{G2-8} از
$!A \lif (!B \lif !C), !A \lif !B \Proves !A \lif !C$ استفاده
می‌کند. کاربرد شرط~P2 در~\olref{G2-5} و~\olref{G2-6} بر نمونه‌هایی از
P2 تکیه دارد:
$\OProv(\gn{!A \lif !B}) \lif (\OProv(\gn{!A}) \lif
\OProv(\gn{!B}))$. در نمونهٔ نخست~$!A \ident !G$ و
$!B \ident \OProv(\gn{!G}) \lif \lfalse$؛ در نمونهٔ دوم
$!A \ident \OProv(\gn{!G})$ و~$!B \ident \lfalse$.

صورت انتزاعی‌تر قضیهٔ دوم ناتمامیت چنین است:

\begin{thm}
\ollabel{thm:second-incompleteness-gen} فرض کنید~$\Th{T}$ نظریه‌ای
دلخواه، سازگار و !!{axiomatized} و گسترشی از~$\Th{Q}$ باشد و
$\OProv[\Th{T}](y)$ هر فرمولی باشد که شرایط !!{derivability} یعنی
P1--P3 را برای~$\Th{T}$ برآورده کند. آنگاه~$\Th{T}$،
$\OCon[\Th{T}]$ را !!{derive} نمی‌کند.
\end{thm}

\begin{prob}
نشان دهید~$\Th{PA}$، $!G_{\Th{PA}} \lif \OCon[\Th{PA}]$ را
!!{derive}s می‌کند.
\end{prob}

\begin{digress}
نتیجهٔ اخلاقی داستان آن است که هیچ نظریهٔ سازگار و «معقول» برای ریاضیات
نمی‌تواند گزارهٔ سازگاری خودش را !!{derive} کند. فرض کنید~$\Th{T}$
نظریه‌ای ریاضی باشد که~$\Th{Q}$ و استدلال «متناهی‌گرایانهٔ» هیلبرت
(هرچه معنایش باشد) را در بر می‌گیرد. آنگاه کل~$\Th{T}$ نمی‌تواند
گزارهٔ سازگاری~$\Th{T}$ را !!{derive} کند؛ پس به طریق اولی، بخش
متناهی‌گرای آن نیز نمی‌تواند گزارهٔ سازگاری~$\Th{T}$ را !!{derive}
کند. به این معنا، اثباتی متناهی‌گرایانه برای سازگاری «تمام ریاضیات»
نمی‌تواند وجود داشته باشد.

در تفسیر اصطلاح «متناهی‌گرایانه» قدری مجال انعطاف وجود دارد و گودل در
مقالهٔ سال ۱۹۳۱ احتمال می‌دهد چیزی که شاید آن را «متناهی‌گرایانه» بدانیم
بیرون از انواع ریاضیاتی قرار گیرد که هیلبرت می‌خواست صوری کند. اما گودل
سخاوتمندانه داوری می‌کرد؛ امروزه دشوار است ببینیم چگونه ممکن است چیزی
بیابیم که بتوان به‌طور معقول متناهی‌گرایانه نامید، اما مثلاً در
$\Th{ZFC}$، یعنی نظریهٔ مجموعه‌های زرملو--فرانکل با اصل انتخاب،
صورت‌بندی‌پذیر نباشد.
\end{digress}

\end{document}
